\section{An exact pullback and the coefficient to be controlled}
\label{sec:generic}

We begin the proof of Theorem~\ref{thm:relative-denominators}.  Put
$k=\dim Z$ and $K=\C(Z)$.  The contraction condition makes $K$
algebraically closed in $\C(X)$, so this function-field extension is
regular in characteristic zero.  The generic fibre is therefore
geometrically integral.  It is normal by localization from $X$, and
normality is geometric for a finite-type variety over the perfect field
$K$.  Choose a rational top form on $Z$.  Together with the form defining
$K_X$, the determinant sequence for function-field differentials determines
a rational relative top form on the generic fibre $X_\eta$.  We use its
divisor as $K_{X_\eta}$.  At codimension-one points of $X_\eta$, which are
smooth over $K$, this divisor is the coefficientwise localization of
$K_X$; vertical primes disappear.  If $d=k$, then $X_\eta=\Spec K$ and
the relative form is a nonzero $0$-form with zero divisor.
Restrict a log resolution and its crepant sub-boundary to the
geometric generic fibre.  Generic smoothness makes the horizontal strata
SNC there, so the original geometric generic pair is lc.  That pair is
projective of dimension $d-k\leq3$, its boundary coefficients belong to
$\Phi$, and its adjoint is rationally linearly trivial.

The lower-dimensional index theorem applies to this pair over
$\overline K$.  For completeness, $K$ is finitely generated over $\C$,
so $\overline K$ and $\C$ are algebraically closed fields of
characteristic zero with the same transcendence degree over $\Q$.
An abstract field isomorphism therefore identifies the pair with a
complex pair to which the theorem applies.

Choose a common integer $p_0$, also clearing $\Phi$, which principalizes
the adjoint of every such geometric generic fibre.  This
principalization descends to $K$ without enlarging the degree.  Indeed,
choose a finite Galois extension $L/K$ over which a principalizing
function is defined.  The quotient of any two of its Galois conjugates
has zero divisor on the proper geometrically integral normal fibre over
$L$, hence belongs to $L^*$.  These quotients form a multiplicative
Galois cocycle.  Hilbert~90 rescales the function to make it invariant,
and it then descends to $K$.  Consequently there is
$\psi\in\C(X)^*$ such that
\[
 p_0(K_X+B)+\Div(\psi)
\]
is vertical over $Z$.

Choose $a$ divisible by $p_0$ and $\phi\in\C(X)^*$ with
\[
 a(K_X+B)+\Div(\phi)=f^*(aD).
\]
The divisor of $\psi^{a/p_0}/\phi$ is vertical.  On the proper normal
generic fibre this function has neither zeros nor poles, so it belongs
to $K^*$.  Denoting it by $h$, set
\[
 D_Z=D+\frac1a\Div(h).
\]
Substitution gives the first equality
of~\eqref{eq:relative-exact-presentation} as an equality of actual
rational divisors.  In particular, $D_Z$ is rational Cartier and
$D_Z\sim_\Q D$.

We can now apply the lc-trivial-fibration theorem
\cite[Theorem~3.6]{FujinoGongyoModuli}.  Besides generic log canonicity
and the rational pullback relation, it requires a rank-one condition.
On a log resolution, the modified discrepancy divisor $\mathbf A^*$
omits the discrepancy $-1$ terms.  Because $0\leq B\leq1$,
rounding up leaves, over the generic point of the base, only an
effective exceptional divisor over $X$: strict transforms contribute
zero.  The associated divisorial sheaf pushes forward to $\OO_X$ by
normality.  Pushing on to $Z$ thus has rank one, as required.

For a prime divisor $P$ on any birational base model, let $t_P$ be the lc
threshold of its pullback over the generic point of $P$, computed with
the crepant sub-pair on a resolved diagram.  The discriminant
coefficient is $1-t_P$.  The crepant sub-pair is sub-lc, so $t_P\geq0$;
in the klt case it is sub-klt and $t_P>0$.  These inequalities on every base model give
the asserted generalized singularities.  On the original base $Z$,
the boundary upstairs is effective and a pullback component
dominating $P$ has positive integral multiplicity.  Hence $t_P\leq1$,
which gives $B_Z\geq0$.

For primes $P$ on $Z$, the thresholds satisfy ACC uniformly.  Shrink
about the generic point of $P$ so that $P$ is Cartier.  In a log
resolution computing the threshold, only divisors with positive order
along the pullback of $P$ matter.  Remove the images of those whose
images are proper subsets of $P$.  Then $t_P$ is an ordinary lc threshold of an effective
Cartier divisor, with positive integral coefficients, against an lc
pair with coefficients in $\Phi$.  The threshold ACC theorem
\cite[Theorem~1.1]{HMX2014} applies.  These thresholds are rational
and lie in $[0,1]$, so their complements form the required rational
DCC set $\mathcal B\subset[0,1]$.

The qualitative theorem makes the moduli b-divisor b-nef and gives its
descent to a nef rational Cartier divisor on a sufficiently high smooth
projective model $W\to Z$.  This holds for our chosen presentation:
changing a rational trivialization adds a rational principal b-divisor
from the base, which preserves nefness and rational Cartier descent.
Write
\[
 D_W=(W\to Z)^*D_Z,
 \qquad M_W=D_W-K_W-B_W.
\]
For $P\subset W$, put $\alpha=\operatorname{coeff}_P D_W$.  Since
$\operatorname{coeff}_P B_W=1-t_P$ and the coefficient of $K_W$ is an
integer,
\[
 \operatorname{coeff}_P M_W
   =\alpha+t_P-1-\operatorname{coeff}_P K_W.
\]
It remains to prove, uniformly in $P$, that
\begin{equation}\label{eq:relative-coefficient-target}
 p(\alpha+t_P)\in\Z.
\end{equation}
Then $pM_W$ is an integral divisor on the smooth variety $W$, hence Cartier.
